\documentclass[12pt,a4paper]{article}
\usepackage[T1]{fontenc}
\usepackage[utf8]{inputenc}
\usepackage{amsfonts,amssymb,amsmath}
\usepackage[dvips]{color}
\usepackage{times}

% D\0x8efinition de la page
\pagestyle{plain}
\usepackage[left=3.5cm, right=3.5cm, top=3cm, bottom=3cm]{geometry}
\newcommand{\vect}[1]{\overrightarrow{\,\mathstrut#1\,}}
\renewcommand{\labelenumii}{\alph{enumii}~)}
\let\le\leqslant
\let\ge\geqslant
% =========================================================================
%                               FIN DU PREAMBULE
% =========================================================================
\begin{document}
% \special{header=ps.pro}
% \special{" cadre1 cadre2}

\begin{center}\Large \color{red}
\sc CONCOURS GÉNÉRAL 1998\\
\rule{10cm}{1mm}
\end{center}

\section*{\color{blue} Exercice 1}

Un tétraèdre ABCD vérifie les conditions suivantes :

\begin{enumerate}
\item les arêtes AB, AC et AD sont deux  à deux orthogonales
\item AB $= 3$ et CD $=\sqrt{2}$.
\end{enumerate}

Déterminer la valeur minimale de $\text{BC}^6 + \text{BD}^6 - \text{AC}^6 - \text{AD}^6$.

\section*{\color{blue} Exercice 2}

Soit $\left(u_n\right)_{n\in \mathbb{N}}$ une suite réelle vérifiant, pour tout
entier $n$, la relation :

\[u_{n+2}= \left|u_{n+1}\right|-u_n.\]

Montrer qu'il existe un entier $p$ non nul,
tel que la relation $u_n= u_{n+p}$ ait lieu pour tout entier naturel~$n$.

\section*{\color{blue} Exercice 3}

Pour tout réel $x$ on note $E(x)$, le plus grand entier
relatif inférieur ou égal à $x$. 

Soit $k$ un entier fixé, supérieur ou égal à 2. 

On considère la fonction $f$ de $\mathbb{N}$ dans $\mathbb{N}$ définie par :

\[f(n)=n+E\big(\root k\of{n+\root k\of{n}}\big)\]

Déterminer l'ensemble des valeurs prises par la fonction $f$.

\section*{\color{blue} Exercice 4}

On considère deux droites $D_1$ et $D_2$ sécantes en $O$,
et un point $M$ n'appartenant à aucune de ces deux droites. On
considère deux points variables, $A$ sur $D_1$ et $B$ sur
$D_2$, tels que le point $M$ appartienne au segment $[A,B]$.

\centerline{(Les questions 1 et 2 sont indépendantes)}

\begin{itemize}
\item[(1)] Montrer qu'il existe une position des points $A$ et $ B$ pour
laquelle l'aire du triangle $OAB$ est minimale.
Construire les points $A$ et $B$ ainsi déterminés.
\item[(2)]  Montrer qu'il existe une position des points
$A$ et $B$ pour laquelle le périmètre du triangle $OAB$ est
minimal et qu'on a alors l'égalité des périmètres des
triangles $OAM$ et $OBM$, ainsi que la relation~:
$$\frac {AM}{\tan  {\widehat{OAM}\over 2}} =
\frac {BM}{\tan  {\widehat{OBM}\over 2}}$$ Construire les point $A$ et
$B$ ainsi déterminés.
\end{itemize}

\section*{\color{blue} Exercice 5}

Soit $n$ un entier supérieur ou égal à 3. On
considère un ensemble $A$ de $n$ points du plan, cet ensemble
ne contenant pas trois points alignés.

Montrer qu'il existe un ensemble $S$ de $2n-5$ points du plan
tel que pour tout triangle dont les sommets sont des points de $A$ il
existe au moins un point de $S$ qui lui soit strictement intérieur.
\end{document}